% ---------------------------------------------------------------------
%  Abstract (150 to 250 words). Write it last.
%  Five moves, one or two sentences each (same order as the proposal):
%   1. Context: offensiveness is subjective; groups and cultures disagree;
%      single labels erase this.
%   2. Problem / gap: LLMs are used as annotators and given personas, but
%      it is unclear whether the persona changes what the model does
%      internally or only its output; prior work lacks content-free
%      controls.
%   3. What we do: D3CODE (8 regions), Llama-3.1-8B-Instruct and
%      Qwen2.5-7B-Instruct; 22 conditions x 10 prompt formats x 500 texts;
%      length-matched placebos and place-name controls; per-head attention;
%      causal head profiling (PASTA); attention and activation steering.
%   4. Findings (current numbers, CLAUDE.md section 9):
%      - only region framing beats same-length placebos (+1.6 / +2.8
%        log-odds); gender and religion do not;
%      - Llama: about half the region effect is identity, half the place
%        name; Qwen (perspective formats): the place name;
%      - framing moves both models AWAY from the real groups
%        (distance 0.35 -> 0.51, 0.40 -> 0.68);
%      - a few middle-layer heads carry the effect causally; steering them
%        halves the distance (0.50 -> 0.22), but a constant shift does as
%        well (0.19); no method creates group-specific alignment;
%      - D3CODE groups differ only by small average offsets, not text by
%        text (split-half agreement about 0).
%   5. Implication: demographic prompting injects a model-made group
%      effect; steering can switch it off but cannot create alignment
%      where the data has no group-specific signal.
%
%  The proposal's draft abstract claimed the placebo pushes the opposite
%  way to real framings. That was a v1 pilot result with one placebo and
%  does not hold in v2 (Qwen placebos push towards "offensive").
% ---------------------------------------------------------------------
\begin{abstract}
\noindent
\fillin{Context.}
\fillin{Problem and gap.}
\fillin{What we do: data, models, design, methods.}
\fillin{Main findings, with two or three numbers.}
\fillin{Implication.}
\end{abstract}
